% !TeX root = code_listings_demo.tex
% !TeX program = pdflatex
% !TeX encoding = UTF-8
\documentclass[t,10pt,aspectratio=169,rounded]{beamer}
\usetheme{CEA23}

\usepackage[utf8x]{inputenc}
\usepackage[english]{babel}
\usepackage{csquotes}

% Code listings: Carbon-style "code snapshot" cards.
\usepackage{beamer-code-listings}

\begin{document}

\title{Code listing demo}
\subtitle{Carbon-style cards for the CEA23 Beamer theme}
\date[31/07/2026]{31 July 2026}
\author{Riccardo \textsc{Finotello}}

\titlepagelogo{\includegraphics[height=24mm]{graphics/ORIGINAL/ISAS.png}}
\setvalue{\diffusion}{}
\seminar[PolytechParis]{Machine Learning - PolytechParis}
\CEAmerci{Thank you}
\CEAcenter{CEA SACLAY}
\addr{91\,191 Gif-sur-Yvette Cedex\\France}
\CEAemail{\href{mailto:riccardo.finotello@cea.fr}{\texttt{riccardo.finotello@cea.fr}}}

\begin{frame}[plain]
	\titlepage
\end{frame}

% --- Feature 1: plain lstlisting with default dark style ---
\begin{frame}[fragile]
	\frametitle{Feature 1: plain \texttt{lstlisting}}
	A plain listing inherits the dark Carbon style.
	\begin{lstlisting}
import pandas as pd

# Load the MARS dataset
df = pd.read_csv("mars.csv")
print(df.head())
print(df.describe())
\end{lstlisting}
\end{frame}

% --- Feature 2: codewithcard dark card ---
\begin{frame}[fragile]
	\frametitle{Feature 2: dark code card}
	\begin{codewithcard}{mars_eda.py}
import pandas as pd
import matplotlib.pyplot as plt

# Load the MARS dataset
df = pd.read_csv("mars.csv")
print(df.head())

# Basic statistics
print(df.describe())
	\end{codewithcard}
\end{frame}

% --- Feature 3: codewithcardlight light card ---
\begin{frame}[fragile]
	\frametitle{Feature 3: light code card}
	\begin{codewithcardlight}{mars_eda_light.py}
import pandas as pd
import matplotlib.pyplot as plt

# Load the MARS dataset
df = pd.read_csv("mars.csv")
print(df.head())

# Basic statistics
print(df.describe())
	\end{codewithcardlight}
\end{frame}

% --- Feature 4: language override (bash / sh) ---
\begin{frame}[fragile]
	\frametitle{Feature 4: language override}
	\begin{codewithcard}[language=bash]{terminal}
# Create a new uv project for the ML course
uv init --python ">=3.10,<3.13" mars-ml
cd mars-ml
uv add pandas numpy scikit-learn
	\end{codewithcard}
\end{frame}

% --- Feature 5: Pattern A (whole-card reveal) ---
% The whole card appears only from slide 2 onwards.
\begin{frame}[fragile]
	\frametitle{Pattern A: whole-card reveal}
	This frame has 2 slides.
	\begin{onlyenv}<2->
		\begin{codewithcard}{train_step_1.py}
import numpy as np
from sklearn.linear_model import LinearRegression

X = np.random.rand(100, 3)
y = X[:, 0] + 2 * X[:, 1] - 0.5 * X[:, 2]
		\end{codewithcard}
	\end{onlyenv}
\end{frame}

% --- Feature 6: Pattern B (line-by-line reveal inside one card) ---
% \pause inside a tcolorbox listing makes the frame title disappear on early
% overlays.  The robust alternative is to wrap several identical-position
% codewithcard environments in onlyenv blocks.  Each card is top-aligned, so the
% reveal looks like a single card growing line-by-line.
\begin{frame}[fragile]{Pattern B: line-by-line reveal inside one card}
	This frame should produce 3 slides.
	\begin{onlyenv}<1>
		\begin{codewithcard}{train_step_2.py}
import numpy as np
from sklearn.linear_model import LinearRegression
		\end{codewithcard}
	\end{onlyenv}
	\begin{onlyenv}<2>
		\begin{codewithcard}{train_step_2.py}
import numpy as np
from sklearn.linear_model import LinearRegression

X = np.random.rand(100, 3)
y = X[:, 0] + 2 * X[:, 1] - 0.5 * X[:, 2]
		\end{codewithcard}
	\end{onlyenv}
	\begin{onlyenv}<3>
		\begin{codewithcard}{train_step_2.py}
import numpy as np
from sklearn.linear_model import LinearRegression

X = np.random.rand(100, 3)
y = X[:, 0] + 2 * X[:, 1] - 0.5 * X[:, 2]

model = LinearRegression()
model.fit(X, y)
print(f"R^2 = {model.score(X, y):.3f}")
		\end{codewithcard}
	\end{onlyenv}
\end{frame}

% --- Feature 7: Pattern C (chunked reveal with plain listings) ---
% Three ceacarbonplain chunks separated by two pauses -> 3 slides.
\begin{frame}[fragile]
	\frametitle{Pattern C: chunked reveal with plain listings}
	This frame should produce 3 slides.
	\begin{lstlisting}[style=ceacarbonplain]
# Chunk 1: imports
import numpy as np
import pandas as pd
\end{lstlisting}
	\pause
	\begin{lstlisting}[style=ceacarbonplain]
# Chunk 2: synthetic data
df = pd.DataFrame({
    "x1": np.random.rand(100),
    "x2": np.random.rand(100),
})
\end{lstlisting}
	\pause
	\begin{lstlisting}[style=ceacarbonplain]
# Chunk 3: first look
print(df.head())
print(df.describe())
\end{lstlisting}
\end{frame}

% Feature 8: notebook input cell
\begin{frame}[fragile]
	\frametitle{Feature 8: notebook input cell}
	\begin{notebookin}
		import numpy as np
		import pandas as pd

		df = pd.read_csv("mars.csv")
		print(df.head())
	\end{notebookin}
\end{frame}

% Feature 9: notebook cell with output
\begin{frame}[fragile]
	\frametitle{Feature 9: notebook cell with output}
	\begin{notebookin}
		print(f"R^2 = {model.score(X, y):.3f}")
	\end{notebookin}
	\begin{notebookout}
		R^2 = 0.932
	\end{notebookout}
\end{frame}

% Feature 10: reveal output after input (Pattern D)
\begin{frame}[fragile]
	\frametitle{Feature 10: reveal output after input}
	\begin{onlyenv}<1->
		\begin{notebookin}
			print(f"R^2 = {model.score(X, y):.3f}")
		\end{notebookin}
	\end{onlyenv}
	\begin{onlyenv}<2->
		\begin{notebookout}
			R^2 = 0.932
		\end{notebookout}
	\end{onlyenv}
\end{frame}

\begin{frame}[plain]
	\usebeamertemplate*{last page}
\end{frame}

\end{document}
